% ===== PRÉAMBULE CANONIQUE — à coller VERBATIM en tête de CHAQUE .tex =====
\documentclass[11pt,a4paper]{article}
\usepackage[utf8]{inputenc}\usepackage[T1]{fontenc}\usepackage{lmodern}
\usepackage[french]{babel}
\usepackage{amsmath,amssymb,mathtools}\usepackage{icomma}
\usepackage{siunitx}\sisetup{locale=FR}  % \qty{}{}, \si{}, \ang{}
\usepackage{xcolor}\usepackage{colortbl}
\usepackage{tikz}\usepackage{pgfplots}\pgfplotsset{compat=1.18}
\usepackage[siunitx,europeanresistors]{circuitikz} % circuits (élec.)
\usepackage[most]{tcolorbox}\usepackage{enumitem}\usepackage{array,booktabs}
\usepackage[a4paper,margin=2.2cm]{geometry}
\usetikzlibrary{arrows.meta,calc,angles,quotes,positioning,patterns,%
  backgrounds,shapes.geometric,decorations.pathmorphing,decorations.markings,intersections}
\definecolor{primary}{RGB}{222,49,99}\definecolor{primarydark}{RGB}{150,22,55}
\definecolor{defcol}{RGB}{0,114,178}\definecolor{methcol}{RGB}{52,92,148}
\definecolor{excol}{RGB}{0,135,90}\definecolor{attncol}{RGB}{200,40,55}
\tcbset{boxbase/.style={enhanced,breakable,boxrule=0.9pt,arc=2.5pt,
  left=3mm,right=3mm,top=2.3mm,bottom=2.3mm,fonttitle=\bfseries,coltitle=white}}
% --- boîtes TUTORIEL (noms EXACTS, ne pas renommer) ---
\newtcolorbox{cle}[1][]{boxbase,colback=defcol!6,colframe=defcol,
  title={\textsc{À retenir}\ifx\relax#1\relax\relax\else\textmd{ : \itshape #1}\fi}}
\newtcolorbox{methode}[1][]{boxbase,colback=methcol!7,colframe=methcol,sharp corners,
  title={\textsc{Méthode}\ifx\relax#1\relax\relax\else\textmd{ --- \itshape #1}\fi}}
\newtcolorbox{exemplebox}[1][]{boxbase,colback=excol!5,colframe=excol,
  title={\textsc{Exemple}\ifx\relax#1\relax\relax\else\textmd{ : \itshape #1}\fi}}
\newtcolorbox{piege}[1][]{boxbase,colback=attncol!6,colframe=attncol,
  title={\textsc{Piège}\ifx\relax#1\relax\relax\else\textmd{ : \itshape #1}\fi}}
\newtcolorbox{remarque}[1][]{boxbase,colback=black!4,colframe=black!45,
  title={\textsc{Remarque}\ifx\relax#1\relax\relax\else\textmd{ : \itshape #1}\fi}}
% --- boîtes EXERCICE / CORRIGÉ (noms EXACTS, auto-comptées) ---
\newtcolorbox[auto counter]{exo}[1][]{boxbase,colback=primary!4,colframe=primary,
  title={\textsc{Exercice}~\thetcbcounter\ifx\relax#1\relax\relax\else\textmd{ --- \itshape #1}\fi}}
\newtcolorbox[auto counter]{corrige}[1][]{boxbase,colback=excol!4,colframe=excol,
  title={\textsc{Corrigé}~\thetcbcounter\ifx\relax#1\relax\relax\else\textmd{ --- \itshape #1}\fi}}
\newcommand{\vv}[1]{\vec{#1}}\newcommand{\ds}{\displaystyle}
\newcommand{\R}{\mathbb{R}}\newcommand{\N}{\mathbb{N}}\newcommand{\Z}{\mathbb{Z}}
% =========================================================================
\begin{document}
% THEME_TITLE: L'énergie d'un satellite
\sisetup{per-mode=symbol}

Un satellite en orbite possède de l'énergie \textbf{cinétique} (il se déplace) et de
l'énergie \textbf{potentielle de gravitation} (il est dans le champ de la Terre). Leur
somme, l'énergie mécanique, est constante sur une orbite circulaire.

\begin{cle}[les trois énergies]
En utilisant $v^2=\dfrac{\mathcal{G}M_T}{R_T+h}$ :
\[
\boxed{\,E_c=\tfrac12 m v^2=\dfrac{\mathcal{G}M_T m}{2(R_T+h)}\,}\qquad
\boxed{\,E_p=-\dfrac{\mathcal{G}M_T m}{R_T+h}\,}
\]
\begin{itemize}[leftmargin=1.5em,topsep=1pt,itemsep=0pt]
  \item $E_c>0$ : énergie cinétique, en joules (\si{\joule}) ;
  \item $E_p<0$ : énergie potentielle (référence nulle à l'infini : un corps
        \emph{lié} a $E_p<0$) ;
  \item $m$ : masse du satellite (\si{\kilogram}) ; \; $r=R_T+h$ : rayon de l'orbite (\si{\metre}).
\end{itemize}
\end{cle}

\begin{cle}[énergie mécanique : $E_m=-E_c$]
\[
E_m=E_c+E_p=\frac{\mathcal{G}M_T m}{2(R_T+h)}-\frac{\mathcal{G}M_T m}{R_T+h}
\quad\Longrightarrow\quad
\boxed{\,E_m=-\dfrac{\mathcal{G}M_T m}{2(R_T+h)}=-E_c\,}
\]
$E_m$ est \textbf{constante} sur l'orbite (circulaire : $v$ et $r$ ne changent pas) et
\textbf{négative} : le satellite est \emph{lié} à la Terre.
\end{cle}

\begin{center}
\begin{tikzpicture}[scale=1.0]
  \draw[->,black!60] (0,-2.0) -- (0,1.4) node[above,font=\small]{énergie};
  \draw[black!30] (-0.15,0)--(6.4,0);
  \node[left,font=\scriptsize] at (0,0) {$0$};
  % Ec > 0
  \fill[excol!70] (0.7,0) rectangle (1.7,1.0);
  \node[font=\scriptsize,above] at (1.2,1.0) {$E_c>0$};
  % Ep < 0 (double hauteur)
  \fill[defcol!70] (2.6,0) rectangle (3.6,-2.0);
  \node[font=\scriptsize,below] at (3.1,-2.0) {$E_p=-2E_c<0$};
  % Em = -Ec
  \fill[attncol!70] (4.5,0) rectangle (5.5,-1.0);
  \node[font=\scriptsize,below] at (5.0,-1.0) {$E_m=-E_c<0$};
\end{tikzpicture}\\[-1pt]
{\footnotesize En valeur : $E_p=-2E_c$ et $E_m=E_c+E_p=-E_c$ (toujours négative).}
\end{center}

\begin{methode}[calculer les énergies]
1) $r=R_T+h$. \quad 2) $E_c=\dfrac{\mathcal{G}M_T m}{2r}$ (ou $\tfrac12 m v^2$). \quad
3) $E_p=-\dfrac{\mathcal{G}M_T m}{r}=-2E_c$. \quad 4) $E_m=E_c+E_p=-E_c$.
\end{methode}

\begin{exemplebox}[satellite de \qty{800}{\kilogram} à \qty{800}{\kilo\metre}]
$r=\qty{7.2e6}{\metre}$, $\mathcal{G}M_T m=\num{4.0e14}\times 800=\num{3.2e17}$ :
\[
E_c=\frac{\num{3.2e17}}{2\times\num{7.2e6}}\approx\qty{2.22e10}{\joule},\quad
E_p=-\frac{\num{3.2e17}}{\num{7.2e6}}\approx-\qty{4.44e10}{\joule},
\]
\[
E_m=E_c+E_p\approx\qty{2.22e10}{}-\qty{4.44e10}{}=-\qty{2.22e10}{\joule}=-E_c.
\]
\end{exemplebox}

\begin{piege}[signes et « énergie de libération »]
$E_p$ et $E_m$ sont \textbf{négatives} (corps lié), $E_c$ positive. Comme $E_m=-E_c$,
pour \emph{libérer} le satellite (l'envoyer à l'infini, $E_m\to 0$) il faut lui
fournir exactement $+E_c$. Ne jamais oublier le signe « $-$ » de $E_p$.
\end{piege}

\end{document}
